\section{Resultados}

\begin{figure}[!ht]
	\centering
	\begin{circuitikz}[american voltages]
		\draw (0,0) to [vR] ++(3,0)
		to [ammeter] ++(2,0)
		coordinate (Vp)
		-- ++(2,0)
		node[circ, label={above:A}] (A) {}
		to [R={$R_x=\qty{0}{\ohm}$}] ++(3,0)
		coordinate (Ra)
		to [R={$R_x=\qty{0}{\ohm}$}] ++(0,-4)
		coordinate (Rb)
		to [R={$R_x=\qty{0}{\ohm}$}] ++(-3,0)
		node[circ, label={below:B}] (B) {}
		-- ++(-2,0)
		coordinate (Vn)
		-- ++(-5,0)
		to [battery1, v<={$\qty{15}{V}$}, invert] (0,0);

		\draw (Ra) to [R={$R_x=\qty{0}{\ohm}$}] ++(3,0)
		node[circ, label={above:C}] (C) {};

		\draw (Rb) -- ++(3,0)
		node[circ, label={below:D}] (D) {};

		\draw (Vp) to [voltmeter] (Vn);


		\draw[dashed, thick, blue] ($(A)+(-0.5,1)$) rectangle ($(D)+(0.5,-1)$);
		\node[blue] at ($(A)!0.5!(C) + (0,1.3)$) {\textbf{testCaption}};
	\end{circuitikz}
	\caption{Montagem experimental}
	\label{circ:montagem1}
\end{figure}






\begin{table}[!ht]
	\centering
	\caption{mont 1 caso a)}
	\label{mont1alinA}
	% Use pgfplotstable to render the CSV
	\pgfplotstabletypeset[
		col sep=comma, % Tells LaTeX the file is a CSV
		string type,   % Treat headers as text
		% Define the full grid (vertical lines)
		column type={|c|},
		% Styling the header and adding the horizontal lines for the "Full Grid"
		every head row/.style={
				before row=\hline,
				after row=\hline
			},
		every last row/.style={
				after row=\hline
			},
		% Assigning column headers and units
		columns/Voltage/.style={
				column name={Tensão (\unit{\volt})},
				column type={|c|},
				% column type={|S[table-format=2.2]|} % Adjust format to your max digits
			},
		columns/Current/.style={
				column name={Corrente (\unit{\milli\ampere})},
				column type={c|},
				% column type={S[table-format=2.1]|} % Adjust format to your max digits
			},
	]{dados/1a.csv}
\end{table}






\begin{table}[!ht]
	\centering
	\caption{mot 1 caso b)}
	\label{tab:mont1alinB}
	% Use pgfplotstable to render the CSV
	\pgfplotstabletypeset[
		col sep=comma, % Tells LaTeX the file is a CSV
		string type,   % Treat headers as text
		% Define the full grid (vertical lines)
		column type={|c|},
		% Styling the header and adding the horizontal lines for the "Full Grid"
		every head row/.style={
				before row=\hline,
				after row=\hline
			},
		every last row/.style={
				after row=\hline
			},
		% Assigning column headers and units
		columns/Voltage/.style={
				column name={Tensão (\unit{\volt})},
				column type={|c|},
				% column type={|S[table-format=2.2]|} % Adjust format to your max digits
			},
		columns/Current/.style={
				column name={Corrente (\unit{\milli\ampere})},
				column type={c|},
				% column type={S[table-format=2.1]|} % Adjust format to your max digits
			},
	]{dados/1b.csv}
\end{table}













\begin{figure}[!ht]
	\centering
	\begin{circuitikz}[american voltages]
		\draw (0,0) to [battery1, v=15V, name=mainbat] ++(0,-4)
		-- ++(2,0)
		node[circ, label={above:B}] (B) {}
		to [R=wtf] ++(2,0) coordinate(placaS)
		to [R=wtf] ++(0,4) coordinate(placaN)
		to [R=wtf] ++(-2,0)
		node[circ, label={above:A}] (A) {}
		-- (0,0);

		\draw (placaN) to [R=test] ++(2,0)
		node[circ, label={above:C}] (C) {}
		-- ++(2,0)
		node[ocirc, label={above:1}] (out1) {};

		\draw (placaS) to [R=test] ++(2,0)
		node[circ, label={above:D}] (D) {}
		-- ++(2,0)
		node[ocirc, label={above:2}] (out2) {};

		\node[draw, dashed, red, thick, inner sep=20,
			fit=(A) (B) (C) (D),
			label={[blue, font=\bfseries]above:testCaption}] (innerRect) {};

		\node[draw, dashed, blue, thick, inner sep=20,
			fit=(mainbat) (innerRect),
			label={[blue, font=\bfseries]above:legendaTeste}] (outerRect) {};
	\end{circuitikz}
	\caption{Montagem experimental}
	\label{circ:fonte-real-sim}
\end{figure}



\begin{figure}[!ht]
	\centering
	\begin{circuitikz}[american voltages]
		\draw (0,0) to [battery1, v=15V, name=mainbat] ++(0,-4)
		-- ++(1,0)
		node[circ, label={above:B}] (B) {}
		to [R, l_=wtf] ++(2,0) coordinate(placaS)
		to [R, l_=wtf] ++(0,4) coordinate(placaN)
		to [R, l=wtf] ++(-2,0)
		node[circ, label={above:A}] (A) {}
		-- (0,0);

		\draw (placaN) to [R, l_=test] ++(2,0)
		node[circ, label={above:C}] (C) {}
		-- ++(1,0)
		node[ocirc, label={above:1}] (out1) {};

		\draw (placaS) to [R, l_=test] ++(2,0)
		node[circ, label={above:D}] (D) {}
		-- ++(1,0)
		node[ocirc, label={above:2}] (out2) {};

		\draw (out1) to [ammeter] ++(2,0)
		coordinate (reostP)
		to [vR] ++(0,-4)
		coordinate (reostN)
		-- (out2);

		\draw (reostP) -- ++(1,0)
		to [voltmeter] ++(0,-4)
		-- (reostN);


		\node[draw, dashed, blue, thick, inner sep=20,
			fit=(mainbat) (A) (B) (C) (D),
			label={[blue, font=\bfseries]above:legendaTeste}] (outerRect) {};
	\end{circuitikz}
	\caption{Montagem experimental}
	\label{circ:montagem2}
\end{figure}



\begin{table}[!ht]
	\centering
	\caption{mot 2}
	\label{tab:mont2}
	% Use pgfplotstable to render the CSV
	\pgfplotstabletypeset[
		col sep=comma, % Tells LaTeX the file is a CSV
		string type,   % Treat headers as text
		% Define the full grid (vertical lines)
		column type={|c|},
		% Styling the header and adding the horizontal lines for the "Full Grid"
		every head row/.style={
				before row=\hline,
				after row=\hline
			},
		every last row/.style={
				after row=\hline
			},
		% Assigning column headers and units
		columns/Voltage/.style={
				column name={Tensão (\unit{\volt})},
				column type={|c|},
				% column type={|S[table-format=2.2]|} % Adjust format to your max digits
			},
		columns/Current/.style={
				column name={Corrente (\unit{\milli\ampere})},
				column type={c|},
				% column type={S[table-format=2.1]|} % Adjust format to your max digits
			},
	]{dados/2.csv}
\end{table}
